\documentclass[a4paper,10pt]{article}
\newcommand\subdir{/home/koen/LaTeX-setup/}
\input{\subdir/LaTeX-files/imports.tex}
\input{\subdir/LaTeX-files/master-internship.tex}
\usepackage{\subdir/LaTeX-files/aastex}
\usepackage{natbib}
\bibliographystyle{\subdir/LaTeX-files/astroadstitle.bst}

%Set the title, Author and date
\title{Notes Week 34}
\author{Koen Essers}
\tutorial{Master Internship}
\date{\today}

%Set the header style
\header{\tutorialstring}

%Begin the document
\begin{document}
\setuplayout
\section{Bug Report}

\begin{minted}[fontsize=\small]{fortran}
!! @param ACMI An array of the masses in AMU of those species.
ACMI(relevant_species) = chem_isos% W(chem_id(j))
A(relevant_species) = chem_isos% Z_plus_N(chem_id(j))


chem_isos% Z_plus_N = chem_isos% Z + chem_isos% N
chem_isos% W = chem_isos% Z_plus_N + chem_isos% mass_excess*(mev_to_ergs/amu/(clight*clight))
\end{minted}

\section{MESA mailing list}

Hi everyone,

I noticed something while inspecting mean molecular weight profiles in MESA and I was wondering if I am observing the expected behaviour.

Skye EOS and FreeEOS appear to use different approaches to calculate the mean molecular weight \(\mu \).
Skye uses integer isotope mass number (\(A_{i} \)), whereas FreeEOS appears to use its own table of elemental atomic weights.

These different approaches cause discrepancies in the mean molecular weight reported by the two different EOS tables.
To show this, I ran the \lstinline[style=output, breaklines=true]|hb_2M| test suite with thermohaline mixing enabled using mesa-r23.05.1.
The MESA reported \(\mu \) changes sharply across the EOS blending region, transitioning between the values obtained using these two approaches.
During the RGB, the transition between the two tables occurs in a radiative region and appears to trigger thermohaline mixing.

The figure below shows the \(\mu \) profile in the RGB-phase of the \lstinline[style=output, breaklines=true]|hb_2M| model, together with \(\mu \) values independently computed with the two approaches assuming complete ionization, and the thermohaline mixing diffusion coefficient.
\begin{figure}[ht]
	\centering
	\incplot{Skye-vs-FreeEOS}
\end{figure}

Is this difference between the two \(\mu \) calculations something that is already known?
And is the resulting molecular weight gradient, and the associated thermohaline mixing expected when blending Skye and FreeEOS?

Thanks!

Best regards,

Koen

\section{Summarizing my Results}

I will first look at what I think is most relevant in each of my weekly notes.

\begin{enumerate}
	\item[\color{keywordscolor}2 -- 3.] This week contains a lot of relevant MESA parameters that I use, and
	      \emph{why}
	      I use them.
	      Other than that, there are not really any relevant results.
	\item[\color{keywordscolor}4.] This week shows my first attempts of RLOF mass transfer with MESA.
	      I also check some dredge-up properties of \citet{rees2024b}-like TPAGB models.
	      There are no
	      relevant \emph{new} results.
	\item[\color{keywordscolor}5.] I test different wind prescriptions for single stars.
	      This is relevant for the methods section of my final thesis, but does not contain \emph{new} results.
	      I also compare different opacity tables.
	      Similarly, this is important for the methods, but does not show something \emph{new}.

	      I also implement the quasi-adiabatic mass loss limit from \citet{temmink2022}.
	      Again something for the methods.
	      This is also the first week where I try some binary models with RLOF during the TPAGB phase of the donor.
	      Not really any interesting results though, since these models still miss wind-mass transfer and tides.
	\item[\color{keywordscolor}6.] Tried to figure out why models with low \(R\) and \(q\) fail.
	      Did not really succeed.
	      Contains a good description of \AE SOPUS tables for methods / introduction.
	\item[\color{keywordscolor}7.] Contains detailed information of the overshooting parameters and comparisons with other methods.
	      Also, MesaGrid class version 1.
	      Not very relevant any more.
	\item[\color{keywordscolor}8.] Contains the first implementation of the tides.
	      This version does still contain some bugs but it is important for the methods section of my thesis.
	      It also contains the first implementation of a model selector, something that I update many times later on.
	      This is also interesting for the methods section of my thesis.
	\item[\color{keywordscolor}9.] Fixes to tides, also the first implementation of different angular momentum losses due to winds, including the \citet{saladino2018} winds.
	      Shows detailed derivations as well for the angular momentum terms.
	\item[\color{keywordscolor}10.] Starts with a nice figure showing the effects of the tides and wind angular momentum prescriptions for the Roche lobe evolution.
	      Contains a lot of validations for the wind and tides, but this is not the complete and (hopefully) bug free setup yet.
	      It does very nicely show how the individual parts of the prescriptions evolve over time.
	      Also notes for the first time the difference between the two wind velocity assumptions and shows how they affect in particular \citet{saladino2018}.
	\item[\color{keywordscolor}11.] Constant wind velocity of 15 km/s for everything in the model.
	      This also contains important "last" tests for the winds + tides prescription.
	      Shows full consistency with \lstinline[style=output, breaklines=true]|evolve.py| code.
	      Also shows some examples of early grid generation code and how I make sure the mesa models use the correct starting values.
	      Pretty important for methods, not so much for results.
	\item[\color{keywordscolor}12.] Lots of \citet{soberman1997} analytic work.
	      Not really any relevant results, but certainly helped with building intuition and understanding.
	      Shows that \citet{soberman1997} predictions of TPAGB-like mass transfer with modest \(\delta \) values
	      can already produce small orbits on the orbit of a couple hundred days.
	      Also shows for the first time explicitly that small differences in initial binary separation can quite drastically change the outcome.
	\item[\color{keywordscolor}13.] Implementation of proper circumbinary torus mass loss and the first grid showing some results.
	      Also the first time encountering the Darwin instability but no nice plots yet.
	      From this point onwards, the basic MESA models do not change a lot anymore, apart from additions that are required for abundances.
	\item[\color{keywordscolor}14.] This week I show in more detail the Darwin instability.
	      This is also the week that I actually enable the fully conservative mass transfer that \citet{temmink2025} talks about.
	\item[\color{keywordscolor}15 -- 16.] Final model saving code.
	      It now saves a single model at the start of a new thermal pulse with information on the maximum radius that was reached in the time before it.
	      This week I also \emph{observed} the sudden mixing that sometimes occurs in the other layers of the star, which are not convective, but I had not yet \emph{understood} what I was seeing.
	      I also look at the
	      eccentricity and find that a lot of the parameter space is \emph{expected} to be circularized if you \emph{assume} \citet{hut1981} style tides.
	      However, even with \citet{hut1981} style tides, there are still some models that can maintain their eccentricity at the start of RLOF.
	\item[\color{keywordscolor}17.] Here I tried to get some semi-analytic bounds for relevant mass-transfer
	      parameters \(\beta \), \(\gamma \) and \(\delta \).
	      It is based on the criterion that the
	      radius of the star cannot ever become larger than the semi-major axis.
	      This is somewhat interesting, \emph{but} I think it is definitely not a major finding.
	\item[\color{keywordscolor}18.] This week I compared the MESA internal tides with my
	      own implementation.
	      They produce slightly different results, but I opt to use
	      my own method because it works well with the wind mass transfer prescription.
	      The MESA tides does not work well with mass losses without having to implement
	      a lot of new custom code.
	\item[\color{keywordscolor}19.] This week I did some more analytic work.
	      I tried to
	      see the effects of \(q, \delta \) and \(\beta \).
	      The most important result is that isotropic re-emission is \emph{much} less important for the change in binary separation than circumbinary disk mass loss.
	      This means isotropic re-emission can be "used" to tune the amount of accreted mass, while
	      \(\delta \) must be tuned to produce realistic final periods.

	      Then I looked into the radial response of mass transfer for different \(\delta \) values.
	      I find that stronger \(\delta \) values can cause a sort of delayed increase in the radius when a lot the mass is lost.

	      In the HR diagram, this looks like a loop going from the right down, then moving to up and left
	      to the Post-AGB constant luminosity track.
	\item[\color{keywordscolor}20.] This week, I first tried to simulate a single star with a lower metallicity, based on metallicity data from observations of barium giants.
	      Later, I change this metallicity to a more appropriate value for barium dwarfs.
	      We see that a lower metallicity slightly alters the evolution of the star, but it is most important for the abundances.

	      I also implement a very important fix to my angular momentum loss prescription to also take into account angular momentum lost from the accretor.

	      I also implement a better stopping condition, making sure the models have finished RLOF.
	      Then I also implement a profile saving method that saves profiles throughout mass loss.
	\item[\color{keywordscolor}21.] Lower metallicity stars have larger dredge-up efficiencies, in line with literature \citep{herwig2005, karakas1902}.
	      Lower metallicity stars also have larger C/O ratios, also in line with literature.

	      I also compare my MESA models with Monash and NuGrid results without success.
	      My later attemps are much better.

	      I show that it is the mass transfer rate that determines how the star responds radially to mass loss.
	      Large mass loss rates can cause the "loop" in the HR diagram, which is a hint at some not
	      super stable mass loss.

	      I also \emph{show} with the MESA models that \(\delta \) is much more important for the final
	      period than \(\beta \), and that a large \(\delta \) leads to larger mass transfer rates as well.

	      Lastly, I look at the metallicity of dwarfs.
	\item[\color{keywordscolor}22.] I show that the amount of dredge-up for low mass (1 to 3\(\; M_\odot \)) increases with mass in the MESA models, just like in the Monash \citet{karakas2016} models.

	      These weeks, I also update the grid generator code and \lstinline[style=output, breaklines=true]|MesaGrid| class to create and analyze the MESA grids more effectively.

	      Here, I also simulate a bulk of the single star models.
	\item[\color{keywordscolor}23.] I use the masses of observed barium dwarfs in combination with
	      the simple binary code to find suitable binary candidates for the creation of barium dwarfs.

	      I find that it is difficult to create barium dwarfs with masses that agree with observations
	      using large donor stars, simply because of the mass ratio.
	      It is also important to
	      interact later on during the TPAGB, because only then has the star had a third dredge-up
	      episode to enrich the envelope.
	\item[\color{keywordscolor}24.] The MESA binary models have \emph{a lot of} convergence problems.
	      I developed a python pipeline that tries to fix these convergence problems, which is very useful.
	\item[\color{keywordscolor}25.] Made sure my post-processed \(M_\text{DUP}\) calculation agrees well with the internal MESA dredge-up calculation which does not track the
	      actual mass.
	      It agrees well.
	\item[\color{keywordscolor}26.]
	\item[\color{keywordscolor}27.]
	\item[\color{keywordscolor}28.]
	\item[\color{keywordscolor}29.]
	\item[\color{keywordscolor}30.]
	\item[\color{keywordscolor}31.]
	\item[\color{keywordscolor}32.]
	\item[\color{keywordscolor}33.]
	\item[\color{keywordscolor}34.]
\end{enumerate}

\section{Conclusions}

\newsection{2.6 Solar mass grid}
\begin{figure}[ht]
	\centering
	\incplot{w34-pcolor-10-3}
\end{figure}

\begin{figure}[ht]
	\centering
	\incplot{w34-pcolor-11-3}
\end{figure}

\begin{figure}[ht]
	\marginfignote{0}{We see a very strong correlation between the [s/Fe] values and the C/O number ratios.
		This is quite different from observations, where this correlation is \emph{much} less strong.
	}
	\centering
	\incplot{w34-s-co}
\end{figure}

\bibliography{master.bib}
\end{document}
